\documentclass[10pt,a4paper]{amsart}
\usepackage{amsmath,amssymb,amsthm,graphicx}
% Retained minimal preamble for the Stacks Project chapter sites-cohomology.tex (Cohomology on Sites).
% Verbatim excerpt of the chapter preamble: Xy-pic setup (native preamble lines 14-18) and the native
% theorem declarations with the mathematical macros (native preamble lines 154-217). Copied mechanically;
% nothing added, removed or redefined. The wrapper class and the amsmath/amssymb/amsthm/graphicx base
% packages are supplied by the assembler, so the class/externaldocument part of the native preamble is omitted.
% For commutative diagrams we use Xy-pic
\usepackage[all]{xy}

% We use 2cell for 2-commutative diagrams.
\xyoption{2cell}
% Theorem environments.
%
\theoremstyle{plain}
\newtheorem{theorem}[subsection]{Theorem}
\newtheorem{proposition}[subsection]{Proposition}
\newtheorem{lemma}[subsection]{Lemma}

\theoremstyle{definition}
\newtheorem{definition}[subsection]{Definition}
\newtheorem{example}[subsection]{Example}
\newtheorem{exercise}[subsection]{Exercise}
\newtheorem{situation}[subsection]{Situation}

\theoremstyle{remark}
\newtheorem{remark}[subsection]{Remark}
\newtheorem{remarks}[subsection]{Remarks}

\numberwithin{equation}{subsection}

% Macros
%
\def\lim{\mathop{\mathrm{lim}}\nolimits}
\def\colim{\mathop{\mathrm{colim}}\nolimits}
\def\Spec{\mathop{\mathrm{Spec}}}
\def\Hom{\mathop{\mathrm{Hom}}\nolimits}
\def\Ext{\mathop{\mathrm{Ext}}\nolimits}
\def\SheafHom{\mathop{\mathcal{H}\!\mathit{om}}\nolimits}
\def\SheafExt{\mathop{\mathcal{E}\!\mathit{xt}}\nolimits}
\def\Sch{\mathit{Sch}}
\def\Mor{\mathop{\mathrm{Mor}}\nolimits}
\def\Ob{\mathop{\mathrm{Ob}}\nolimits}
\def\Sh{\mathop{\mathit{Sh}}\nolimits}
\def\NL{\mathop{N\!L}\nolimits}
\def\CH{\mathop{\mathrm{CH}}\nolimits}
\def\proetale{{pro\text{-}\acute{e}tale}}
\def\etale{{\acute{e}tale}}
\def\QCoh{\mathit{QCoh}}
\def\Ker{\mathop{\mathrm{Ker}}}
\def\Im{\mathop{\mathrm{Im}}}
\def\Coker{\mathop{\mathrm{Coker}}}
\def\Coim{\mathop{\mathrm{Coim}}}

% Boxtimes
%
\DeclareMathSymbol{\boxtimes}{\mathbin}{AMSa}{"02}

%
% Macros for moduli stacks/spaces
%
\def\QCohstack{\mathcal{QC}\!\mathit{oh}}
\def\Cohstack{\mathcal{C}\!\mathit{oh}}
\def\Spacesstack{\mathcal{S}\!\mathit{paces}}
\def\Quotfunctor{\mathrm{Quot}}
\def\Hilbfunctor{\mathrm{Hilb}}
\def\Curvesstack{\mathcal{C}\!\mathit{urves}}
\def\Polarizedstack{\mathcal{P}\!\mathit{olarized}}
\def\Complexesstack{\mathcal{C}\!\mathit{omplexes}}
% \Pic is the operator that assigns to X its picard group, usage \Pic(X)
% \Picardstack_{X/B} denotes the Picard stack of X over B
% \Picardfunctor_{X/B} denotes the Picard functor of X over B
\def\Pic{\mathop{\mathrm{Pic}}\nolimits}
\def\Picardstack{\mathcal{P}\!\mathit{ic}}
\def\Picardfunctor{\mathrm{Pic}}
\def\Deformationcategory{\mathcal{D}\!\mathit{ef}}

% --- appended 2026-10-12 (Chief): the project own definitions of the reference/slogan/history
% display environments, copied mechanically from the shared Stacks preamble lines 8-12
% (\usepackage{verbatim} plus the three \newenvironment lines, which the chapter preamble keeps
% verbatim from the project preamble). Nothing added, removed or redefined: these environments
% carry no mathematical content and the project itself typesets them as comments. Without these
% three lines a counted statement whose own text uses them fails with Environment slogan undefined.
\usepackage{verbatim}
\newenvironment{reference}{\comment}{\endcomment}
\newenvironment{slogan}{\comment}{\endcomment}
\newenvironment{history}{\comment}{\endcomment}
\usepackage{cleveref}
\crefname{originallemma}{original lemma}{original lemmas}
\Crefname{originallemma}{Original lemma}{Original lemmas}
\makeatletter
\newlabel{lemma-is-limit}{{08U3}{source}{}{sourcepointer-1}{}}
\newlabel{lemma-is-limit@cref}{{[originallemma][0][]08U3}{[1][0][]source}{}{}{}}
\makeatother
\crefname{originallemma}{original lemma}{original lemmas}
\Crefname{originallemma}{Original lemma}{Original lemmas}
\makeatletter
\newlabel{derived-lemma-special-inverse-system}{{070F}{source}{}{sourcepointer-2}{}}
\newlabel{derived-lemma-special-inverse-system@cref}{{[originallemma][0][]070F}{[1][0][]source}{}{}{}}
\makeatother
\crefname{originallemma}{original lemma}{original lemmas}
\Crefname{originallemma}{Original lemma}{Original lemmas}
\makeatletter
\newlabel{derived-lemma-bounded-below-injectives-K-injective}{{070J}{source}{}{sourcepointer-3}{}}
\newlabel{derived-lemma-bounded-below-injectives-K-injective@cref}{{[originallemma][0][]070J}{[1][0][]source}{}{}{}}
\makeatother
\crefname{originallemma}{original lemma}{original lemmas}
\Crefname{originallemma}{Original lemma}{Original lemmas}
\makeatletter
\newlabel{derived-lemma-limit-K-injectives}{{070L}{source}{}{sourcepointer-4}{}}
\newlabel{derived-lemma-limit-K-injectives@cref}{{[originallemma][0][]070L}{[1][0][]source}{}{}{}}
\makeatother
\crefname{originallemma}{original lemma}{original lemmas}
\Crefname{originallemma}{Original lemma}{Original lemmas}
\makeatletter
\newlabel{derived-lemma-difficulty-K-injectives}{{070M}{source}{}{sourcepointer-5}{}}
\newlabel{derived-lemma-difficulty-K-injectives@cref}{{[originallemma][0][]070M}{[1][0][]source}{}{}{}}
\makeatother
\crefname{originallemma}{original lemma}{original lemmas}
\Crefname{originallemma}{Original lemma}{Original lemmas}
\makeatletter
\newlabel{lemma-is-limit-dimension}{{0D6P}{source}{}{sourcepointer-6}{}}
\newlabel{lemma-is-limit-dimension@cref}{{[originallemma][0][]0D6P}{[1][0][]source}{}{}{}}
\makeatother
\title{A sufficient condition for the canonical map into the candidate K-injective complex to be a quasi-isomorphism}
\author{The Stacks Project Authors}
\date{stacks-project:sites-cohomology.tex:local-snapshot-20261011; GNU FDL 1.2 or later}
\begin{document}
\maketitle
\noindent\textit{Editorial note.} All delivered mathematical text is the author's own, copied from the pinned Stacks Project chapter file sites-cohomology.tex; no statement or proof step was rewritten, reordered or removed.\par\medskip
\noindent\textit{Editorial note.} The proof is the author's own argument: it names the general criterion reducing the assertion to the derived limit statement and the companion criterion which supplies it. Both sentences are kept verbatim inside the counted proof span.\par\medskip
\noindent\textit{Editorial note.} Every labelled reference occurring inside the counted spans is carried in assembly.external\_references as a non-load-bearing navigation pointer (same-chapter ones at their definition site, cross-chapter ones at the invocation inside the counted span), so no load-bearing mathematics is replaced by a stub. No bibliography citation occurs inside the counted spans.\par\medskip
\noindent\textit{Editorial note.} The checked chapter list for same-name collisions is the set of stacks chapters staged in this repository: chow, cohomology, divisors, etale-cohomology, flat, more-algebra, more-morphisms, morphisms, spaces-more-morphisms, varieties together with this card's chapter sites-cohomology.\par\medskip
\noindent\textit{Editorial note.} The preamble is the chapter's own retained minimal preamble (sources/stacks-sites-cohomology/retained-preamble.tex), a verbatim excerpt of the native preamble: the Xy-pic setup and the native theorem declarations with the mathematical macros.\par\medskip
\noindent\textit{Allowed standard prerequisites (editorial).}\begin{itemize}
\item The standing set-up: the construction of the candidate K-injective complex and the canonical map into it. Same chapter, the set-up of the section delivered verbatim as a context span in this card (source lines 5548-5584), including the displayed equation that the counted statement refers to. Supplies the map whose being a quasi-isomorphism is asserted.
\item The companion criterion recovering an object as the derived limit of its truncations, and the general K-injectivity criterion. Same chapter, the companion results on the derived limit of the truncations and on K-injective complexes. Supplies the two results to which the proof reduces the assertion.
\end{itemize}\medskip
\ifdefined\hypertarget\hypertarget{sourcepointer-1}{}\fi
\noindent\textit{Source reference.} Original lemma 08U3: sources/stacks-sites-cohomology/pinned-original-source.tex, source lines 5417--5417. Navigation pointer invoked at proof line 5598; the counted span contains the pointer only and the cited result is not reproduced inside it, so the pointer is registered as a non-load-bearing navigation stub in the established form for Stacks material. No bibliography citation is involved.\par\medskip
\ifdefined\hypertarget\hypertarget{sourcepointer-2}{}\fi
\noindent\textit{Source reference.} Original lemma 070F: sources/stacks-sites-cohomology/pinned-original-source.tex, source lines 5554--5554. Navigation pointer invoked at proof line 5554; the counted span contains the pointer only and the cited result is not reproduced inside it, so the pointer is registered as a non-load-bearing navigation stub in the established form for Stacks material. No bibliography citation is involved.\par\medskip
\ifdefined\hypertarget\hypertarget{sourcepointer-3}{}\fi
\noindent\textit{Source reference.} Original lemma 070J: sources/stacks-sites-cohomology/pinned-original-source.tex, source lines 5575--5575. Navigation pointer invoked at proof line 5575; the counted span contains the pointer only and the cited result is not reproduced inside it, so the pointer is registered as a non-load-bearing navigation stub in the established form for Stacks material. No bibliography citation is involved.\par\medskip
\ifdefined\hypertarget\hypertarget{sourcepointer-4}{}\fi
\noindent\textit{Source reference.} Original lemma 070L: sources/stacks-sites-cohomology/pinned-original-source.tex, source lines 5576--5576. Navigation pointer invoked at proof line 5576; the counted span contains the pointer only and the cited result is not reproduced inside it, so the pointer is registered as a non-load-bearing navigation stub in the established form for Stacks material. No bibliography citation is involved.\par\medskip
\ifdefined\hypertarget\hypertarget{sourcepointer-5}{}\fi
\noindent\textit{Source reference.} Original lemma 070M: sources/stacks-sites-cohomology/pinned-original-source.tex, source lines 5604--5604. Navigation pointer invoked at proof line 5604; the counted span contains the pointer only and the cited result is not reproduced inside it, so the pointer is registered as a non-load-bearing navigation stub in the established form for Stacks material. No bibliography citation is involved.\par\medskip
\ifdefined\hypertarget\hypertarget{sourcepointer-6}{}\fi
\noindent\textit{Source reference.} Original lemma 0D6P: sources/stacks-sites-cohomology/pinned-original-source.tex, source lines 5440--5440. Navigation pointer invoked at proof line 5607; the counted span contains the pointer only and the cited result is not reproduced inside it, so the pointer is registered as a non-load-bearing navigation stub in the established form for Stacks material. No bibliography citation is involved.\par\medskip

\noindent
Let $(\mathcal{C}, \mathcal{O})$ be a ringed site.
Let $\mathcal{F}^\bullet$ be a complex of $\mathcal{O}$-modules.
The category $\textit{Mod}(\mathcal{O})$ has enough injectives, hence
we can use
Derived Categories, Lemma \ref{derived-lemma-special-inverse-system}
produce a diagram
$$
\xymatrix{
\ldots \ar[r] &
\tau_{\geq -2}\mathcal{F}^\bullet \ar[r] \ar[d] &
\tau_{\geq -1}\mathcal{F}^\bullet \ar[d] \\
\ldots \ar[r] & \mathcal{I}_2^\bullet \ar[r] & \mathcal{I}_1^\bullet
}
$$
in the category of complexes of $\mathcal{O}$-modules such that
\begin{enumerate}
\item the vertical arrows are quasi-isomorphisms,
\item $\mathcal{I}_n^\bullet$ is a bounded below complex of injectives,
\item the arrows $\mathcal{I}_{n + 1}^\bullet \to \mathcal{I}_n^\bullet$
are termwise split surjections.
\end{enumerate}
The category of $\mathcal{O}$-modules has limits (they are computed
on the level of presheaves), hence we can form the termwise limit
$\mathcal{I}^\bullet = \lim_n \mathcal{I}_n^\bullet$. By
Derived Categories, Lemmas
\ref{derived-lemma-bounded-below-injectives-K-injective} and
\ref{derived-lemma-limit-K-injectives}
this is a K-injective complex. In general the canonical map
\begin{equation}
\label{equation-into-candidate-K-injective}
\mathcal{F}^\bullet \to \mathcal{I}^\bullet
\end{equation}
may not be a quasi-isomorphism. In the following lemma we describe some
conditions under which it is.


\begin{lemma}
\label{lemma-K-injective}
In the situation described above. Denote
$\mathcal{H}^m = H^m(\mathcal{F}^\bullet)$ the $m$th cohomology sheaf.
Let $\mathcal{B} \subset \Ob(\mathcal{C})$ be a subset.
Let $d \in \mathbf{N}$.
Assume
\begin{enumerate}
\item every object of $\mathcal{C}$ has a covering whose members are
elements of $\mathcal{B}$,
\item for every $U \in \mathcal{B}$ we have $H^p(U, \mathcal{H}^q) = 0$
for $p > d$ and $q < 0$\footnote{It suffices if
$\forall m$, $\exists p(m)$, $H^p(U. \mathcal{H}^{m - p}) = 0$ for
$p > p(m)$, see Lemma \ref{lemma-is-limit}.}.
\end{enumerate}
Then (\ref{equation-into-candidate-K-injective}) is a quasi-isomorphism.
\end{lemma}

\begin{proof}
By Derived Categories, Lemma \ref{derived-lemma-difficulty-K-injectives}
it suffices to show that the map
$\mathcal{F}^\bullet \to R\lim \tau_{\geq -n} \mathcal{F}^\bullet$
is an isomorphism. This follows from Lemma \ref{lemma-is-limit-dimension}.
\end{proof}

\end{document}
